\section{Casos y operación: implementación de GR00T fuera de Berkeley}

\subsection{Resumen de casos de la industria}

Jim Fan analizó dos casos reales:

\begin{itemize}
    \item \textbf{Laboratorio de obleas de semiconductores}: usa el modelo GR00T junto con un brazo robótico de sujeción autónomo para completar el intercambio de wafer. Se apoya en built-in alignment rules para evitar colisiones.
    \item \textbf{Centro de clasificación logística}: el objetivo es recoger objetos de formas diversas en un entorno abarrotado, generando planes de ejecución mediante instrucciones de LLM y detectando errores mediante per-instance vision.
\end{itemize}

\begin{warningbox}{La incomodidad real al implementar}
En los escenarios reales, la falla más común no es un error de inferencia del modelo, sino sensor calibration drift, el backlash mecánico y las pequeñas desviaciones que provoca el cambio de herramienta. Incluso el modelo más potente necesita un fallback confiable.
\end{warningbox}

\subsection{Métricas de operación y monitoreo}

Para garantizar la implementación, la cadena de observación se divide en tres capas:

\begin{itemize}
    \item \textbf{Trace}: cada acción/instrucción incluye timestamp, la secuencia de tokens y el tool call log
    \item \textbf{Metrics}: tasa de éxito, cycle time, safety intervention frequency
    \item \textbf{Alerts}: los thresholds preestablecidos (colisión, tiempo de espera agotado de la tarea, planner oscillation) activan la luz roja al dispararse
\end{itemize}

\begin{knowledgebox}{Trace-driven postmortem}
Guardar el trace buffer permite reproducir rápidamente lo ocurrido después de un incident: cada token, cada llamada a herramienta y cada transición de estado puede rastrearse, lo que ayuda a los ingenieros a localizar «qué instrucción del paso 37 provocó la desviación».
\end{knowledgebox}

\subsection{Flujo operativo y Playbook}

Cada proyecto de manipulación sigue un Operational Playbook de cinco pasos:

\begin{enumerate}
    \item \textbf{Plan}: definir la tarea, configurar el documento CTA y generar failure hypotheses
    \item \textbf{Deploy}: escribir la live config, calibrar los sensors e insertar el override switch
    \item \textbf{Observe}: monitorear drop rates, collisions y latency mediante metrics + open telemetry
    \item \textbf{Adjust}: si se dispara una alert, entrar en modo «pause + analyze», reproducir el trace y ajustar la instruction
    \item \textbf{Document}: escribir un postmortem por cada incident y actualizar el Checklist
\end{enumerate}

\begin{warningbox}{Pause + Analyze es el freno más seguro}
Cuando el desempeño se degrada, no ajuste el large language predicate a ciegas; basta con seguir el checklist para poner en pausa el robot, reproducir el trace y revisar el lockstep metric para evitar un cascade failure.
\end{warningbox}

\subsection{Resumen de la sección}

La estabilidad operativa requiere el respaldo de Case study, monitoreo/trace/backstop, y debe combinarse con el playbook para poder enfrentar los incident de la implementación real.

\newpage

